% Generated draft: requires mathematical and source review.
\documentclass[11pt]{article}
\usepackage[margin=25mm]{geometry}
\usepackage{fontspec}
\setmainfont{Latin Modern Roman}
\usepackage{amsmath,amssymb,mathtools,unicode-math}
\setmathfont{Latin Modern Math}
\usepackage{enumitem,needspace,xurl,hyperref}
\setlength{\parindent}{0pt}
\setlength{\parskip}{5pt}
\setlength{\emergencystretch}{3em}
\newcommand{\N}{\mathbb{N}}
\newcommand{\Z}{\mathbb{Z}}
\newcommand{\Q}{\mathbb{Q}}
\newcommand{\R}{\mathbb{R}}
\newcommand{\C}{\mathbb{C}}
\newcommand{\F}{\mathbb{F}}
\newcommand{\A}{\mathbb{A}}
\newcommand{\PP}{\mathbb{P}}
\newcommand{\E}{\mathbb{E}}
\newcommand{\eps}{\varepsilon}
\newcommand{\dd}{\,\mathrm{d}}
\begin{document}
\section*{UnsolvedMath ID: 800490}
\textbf{Source problem code:} catalogue entry
\subsection{Definitions and mathematical statement}
For a positive integer \(n\), the standard symplectic torus is conventionally
\[
(\mathbb T^{2n},\omega_0)=\left(\R^{2n}/\Z^{2n},\sum_{j=1}^{n}\dd x_j\wedge\dd y_j\right).
\]
For a closed symplectic manifold \((M,\omega)\), let \(\operatorname{Ham}(M,\omega)\) denote the Hamiltonian diffeomorphisms: time-one maps of flows generated by smooth time-dependent Hamiltonians on \(M\). Here closed means compact without boundary.

The supplied question asks whether every \(\varphi\in\operatorname{Ham}(\mathbb T^{2n},\omega_0)\) satisfies
\[
\liminf_{T\to+\infty}\frac{\log P_\varphi(T)}{\log T}\geq 2,
\]
and whether the same assertion holds for every Hamiltonian diffeomorphism of an arbitrary aspherical closed symplectic manifold. The notation \(P_\varphi(T)\) merely makes the presumed dependence of the source's \(P(T)\) on \(\varphi\) explicit; it does not supply a definition.

\textbf{Unresolved qualifications.} The complete supplied record does not define \(P(T)\). In particular, it does not identify what is counted, any equivalence relation used in counting, the meaning of the parameter \(T\), or whether \(T\) ranges over integers or real numbers. The displayed logarithm requires a positive finite value of \(P_\varphi(T)\), unless some extended-value convention is specified; none is supplied. Thus the displayed question cannot yet be made into a fully specified mathematical assertion. The torus formulation above uses the conventional standard torus; the source does not specify its dimension or explicitly state the quantification over dimensions.

The meaning of aspherical is also unspecified. Possible hypotheses include \(\pi_2(M)=0\), vanishing of \(\int_{S^2}u^*\omega\) for every smooth \(u:S^2\to M\), or the latter condition together with \(\langle c_1(TM),u_*[S^2]\rangle=0\) for every such \(u\), where \(c_1(TM)\) is computed using an \(\omega\)-compatible almost complex structure. These hypotheses must not be silently identified; the intended one requires clarification from the source.
\subsection{Short English statement}
The question asks whether the quantity denoted by \(P(T)\) has lower logarithmic growth exponent at least \(2\) for every Hamiltonian diffeomorphism of a standard symplectic torus, and more generally of an aspherical closed symplectic manifold. The supplied record leaves \(P(T)\) undefined and does not specify the asphericity hypothesis, so this is an incomplete statement-only draft requiring mathematical review.
\subsection{Sources}
\begin{itemize}
\item Supplied local UnsolvedMath registry record: permanent ID \(800490\), exact source code \texttt{catalogue entry}, upstream revision \texttt{d371ee917ac9647fbb06cd6fc529deec6f3dda82}. The registry labels the problem open; this is supplied curation, not an independently verified mathematical or literature assessment.
\item Supplied source attribution: Albach, Barraud et al., \emph{Open Problems in Billiards and Quantitative Symplectic Geometry} (2026), \url{https://arxiv.org/abs/2602.12896}. The record attributes the item and contribution to Bernhard Albach, Umberto Hryniewicz, and Alexandru Oancea. It records access on 2026-10-16 and extraction from arXiv LaTeX; these are supplied provenance notes, not a claim of independent access or verification.
\item The background and title identify the source item as Problem 3, while the extracted statement begins ``Problem 2.'' The supplied evidence does not resolve this discrepancy.
\end{itemize}
\end{document}
